Il seguente Capitolo fornirà un'introduzione teorica al concetto di esecuzione non sequenziale di un processo, ponendo enfasi sui vantaggi portati e sulle nuove necessità richieste da tale implementazione. 

\section[Algoritmo, Programma, Processo]{Algoritmo, Programma, Processo}
Nella definizione della singola unità di esecuzione si ritiene opportuno effettuare una leggera digressione. \par
Un Algoritmo è un procedimento puramente logico studiato appositamente per la risoluzione di uno specifico problema. Nella letteratura sono rintracciabili moltissimi Algoritmi, per brevità vengono nominati a titolo di esempio Algoritmi di Ordinamento come il \textit{Bubble Sort}\cite{BubbleSort} od il \textit{Radix Sort}\cite{RadixSort}, e Algoritmi di Ricerca del Cammino Minimo in un Grafo come l'\textit{Algoritmo di Dijkstra}\cite{DijksraAlg} o l'\textit{Algoritmo di Bellman-Ford}\cite{BellmanFordAlg}. \par
Un Programma è invece la descrizione in codice di un Algoritmo in uno specifico Linguaggio di Programmazione in maniera tale da permetterne l'esecuzione all'interno di un Calcolatore. Si riporta di seguito un'esempio di Programma che descriva l'Algoritmo di Ordinamento per confronti ``\textit{QuickSort}''\cite{QuickSort} implementato tramite il Linguaggio di Programmazione Java.
% TODO: QuickSort in Java e Descrizione di Processo.